% Page 047

\seite{47}

\abschnitt{Pfarrwechsel}
\pstart
Am 11.\ des\unsicher{} ging auf seinen Wunsch Herr Pfarrer Schmitt aus\ob
der Pfarrei. Um seiner Versetzung willen hielt er zum\ob
letzten Male Gottesdienst. Auf die Pfarrei zu Seffern zum\ob
Nachfolger wurde Herr Peter Zimmer, bis dahin Kaplan in\ob
Befort\unsicher, ernannt, welcher durch den Herrn Dechanten\ob
feierlich am 28.\ Mai\unsicher{} eingeführt wurde.

\pend
\abschnitt{Schulfeierlichkeiten}
\pstart

Am Geburtstage des\unsicher{} hundertjährigen\unsicher{} Geburtstages der aller\ohb
gnädigsten\unsicher{} wurde wie das Mal von allen Schulen der Pfarrei\ob
Seffern gemeinschaftlich mit Patrioten\unsicher{} in Seffern\unsicher\ob
gefeiert. Abends gab Herr Götz ein Bier\unsicher.

\pend
\abschnitt{Wechsel in der Ortsschulinspektion}
\pstart

An Stelle des nach\unsicher{} versetzten Herrn Pfarrer\ob
Schmitt wurde am 28.\ Juli durch Verfügung Herr Pfarrer\ob
Zimmer zu Seffern zum Ortsschulinspektor ernannt.

\pend
\abschnitt{Spaziergang}
\pstart

Am 6.\ September unternahmen die vier Schulen der Pfarrei ge\ohb
meinschaftlich einen Spaziergang nach Meckelhausen und\ob
Hamm.

\pend
\abschnitt{Ernte}
\pstart

\margmark{Seffern\\Sefferweich\\15. Dezember}%
 Die diesjährige Ernte kann als eine befriedigende angesehen\ob
werden. Wegen vieler Regengüsse\unsicher{} wurde die Heu- und\ob
Kleeernte in die Länge, alles übrige Getreide\ob
wurde bei günstiger Witterung gut eingebracht. Obst gab es\ob
\margmark{Zimmer\\Dimmer}%
stellenweise viel\unsicher. Kartoffeln gab es wenig\unsicher. Die Viehpreise\ob
sind hoch.
\pend
